% !PRINTED none
\begin{center}
\gesperrt{Neunzehnter Abschnitt.}

\gesperrt{Beziehungen eines K\"orpers auf seine Theiler.}

\rule{0.2\linewidth}{0.4pt}

\S.\,159.

\gesperrt{Relativnormen.}
\end{center}

Wir wollen in diesem Abschnitte die Modificationen betrachten, die in den Definitionen und S\"atzen der Theorie der algebraischen Zahlen eintreten, wenn an Stelle des \gesperrt{absoluten Rationalit\"atsbereiches}, d. h. des K\"orpers der rationalen Zahlen, ein beliebiger algebraischer Zahlk\"orper gesetzt wird. Es werden sich dabei einige wichtige Erg\"anzungen auch f\"ur die allgemeine Theorie der Primfactoren ergeben, die in einer Reihe von \gesperrt{Dedekind} und \gesperrt{Hilbert} aufgestellter S\"atze ihren Ausdruck finden${}^{1}$).

Es sei also $R$ ein algebraischer Zahlk\"orper $m^{\text{ten}}$ Grades und
\weq{1}{f(\Theta) = \Theta^n + \alpha_1\Theta^{n-1} + \cdots + \alpha_n = 0}
eine in $R$ rationale und irreducible Gleichung $n^{\text{ten}}$ Grades. Der K\"orper $\Omega = R(\Theta)$ ist ein algebraischer K\"orper \gesperrt{\"uber} $R$, der in Bezug auf $R$ vom $n^{\text{ten}}$ Grade ist (Bd. I, \S.\,142). Im absoluten Rationalit\"atsbereiche ist $\Omega$ vom Grade $mn$, und wenn $\varrho$ eine Primitivzahl des K\"orpers $R$ ist, so ist
\weq{2}{\Theta^s \varrho^t, \qquad \begin{aligned} s &= 0, 1, \ldots, n-1\\ t &= 0, 1, \ldots, m-1 \end{aligned}}
eine Basis des K\"orpers $\Omega$; denn wegen der Irreducibilit\"at der
%CONT>>

${}^{1}$) \gesperrt{Dedekind}, \glqq Ueber die Discriminanten endlicher K\"orper\grqq{}, Abhandlungen der G\"ottinger Gesellschaft der Wissenschaften (1882). \glqq Zur Theorie der Ideale\grqq{}, Nachrichten der Gesellschaft der Wissenschaften zu G\"ottingen, 1894, Nr. 4. \gesperrt{Hilbert}, \glqq Grundz\"uge einer Theorie des \gesperrt{Galois'schen} Zahlk\"orpers\grqq{}. Ebend. 1894, Nr. 3.
